\section{Seized for the Altar: Presbyter and Bishop of Hippo}

\subsection{The ordination of 391}

In 391 Augustine came to Hippo Regius, the port city two hundred miles from
Thagaste, to recruit a friend for the community---safe, he thought, because
Hippo had a bishop. But the aged Valerius happened to preach that day on the
need for a presbyter, and the congregation knew who was standing among them.
``The Catholics, already acquainted with the life and teaching of the holy
Augustine, laid hands on him---for he was standing there among the people
secure and unaware of what was about to happen\ldots\ and, as is the custom
in such cases, brought him to the bishop to be ordained, for all with common
consent desired that this should be done\ldots\ and they demanded it with
great zeal and clamor, while he wept freely'' (Possidius, \work{Vita} 4).
Bystanders, Possidius adds, took the tears for wounded pride---surely the
famous convert wanted more than a presbyterate---and consoled him that the
office ``was but little inferior to the bishopric''; Augustine himself later
said he wept because he ``apprehended the many imminent dangers'' of
governing a church (\work{Vita} 4). Forced vocation of this kind was common
in the African church; what was uncommon was what the conscript did with it.
His first letter as a cleric begs Valerius for a working delay: knowing his
``infirmity,'' his duty is ``to study with diligence all the remedies which
the Scriptures contain for such a case as mine, and to make it my business
by prayer and reading'' to become fit ``for labours so responsible''
(\work{Letter} 21, of 391). And he asked for his community: ``Soon after he
had been made presbyter he established a
monastery within the church and began to live with the servants of God
according to the manner and rule instituted by the holy apostles. The
principal rule of this society was that no one should possess anything of
his own, but that all things should be held in common''
(\work{Vita} 5). From that monastery, once its founder was a bishop, the
African churches ``began to demand with great eagerness bishops and
clergy''---``about ten men, holy and venerable, continent and learned,''
Possidius counts, himself among them, who in turn founded monasteries of
their own (\work{Vita} 11).

Valerius, ``a Greek by birth and less versed in the Latin language and
literature,'' then broke African custom deliberately: he ``gave his
presbyter the right of preaching the Gospel in his presence in the church
and very frequently of holding public discussions---contrary to the practice
and custom of the African churches. On this account some bishops found fault
with him'' (\work{Vita} 7). The presbyter's public debut in that role
included a public disputation with the Manichaean presbyter Fortunatus,
whose teaching had ``infected and permeated very many'' at Hippo;
Fortunatus, unable to answer, ``left the city of Hippo soon after and
returned to it no more'' (\work{Vita} 6).

\subsection{Coadjutor bishop, 395/396}

Fearing that some vacant see would steal his presbyter, Valerius arranged
with the primate of Carthage to have Augustine consecrated as his own
coadjutor. When Megalius of Calama, primate of Numidia, came to Hippo,
the people acclaimed the plan; Augustine himself ``refused to accept the
episcopate contrary to the custom of the Church, since his bishop was still
living,'' and yielded only when shown precedents ``from the churches across
the sea as well as in Africa'' (Possidius, \work{Vita} 8). Possidius then
records, with a candor that marks his whole book, that Augustine came to
regret the arrangement and said so publicly: it should not have been done
``because of the prohibition of the Ecumenical Council of which he learned
after his ordination\ldots\ for that which he regretted to have had done in
his case he did not wish to have done to others'' (\work{Vita} 8). Thirty
years later, designating his own successor, Augustine told his people
plainly that he ``did not know'' the double episcopate ``to have been
forbidden by the Council of Nicaea'' and refused to repeat the mistake:
Eraclius would remain a presbyter while Augustine lived (\work{Letter} 213,
the recorded acta of 26 September 426). Valerius died soon
after, and Augustine held the see of Hippo for some thirty-four years.

\witnesslimit{The year of the consecration is not directly attested.
Official reception states it as 395 (Benedict XVI, General Audience, 9
January 2008: ``Bishop from 395 to his death in 430''); the older reference
chronology derives 396 from the Encyclopedia's report that ``he was then
forty two''; modern accounts print 395 or 396. The event and its
irregularity are secure on Possidius's witness; this study writes 395/396
and decides nothing further.}

\subsection{The working bishop}

Possidius's central chapters preserve, from forty years' observation, the
texture of an African episcopate: the daily arbitration of lawsuits between
strangers, which Augustine ``regarded as a kind of conscription''
(\work{Vita} 19); intercession with the civil power for prisoners
(\work{Vita} 20); garments, foot-wear, and bedclothing ``modest yet
sufficient---neither too fine nor yet too mean''; a table ``frugal and
sparing,'' herbs and lentils with meat for guests and the weak, wine always,
silver spoons and earthen dishes (\work{Vita} 22); church business held at
arm's length, with no appetite for new buildings and no acceptance of
legacies that stripped a dead man's children (\work{Vita} 23--24); and a
household rule stricter than the canons required: ``No woman ever lived or
stayed in Augustine's house, not even his own sister,'' a widow who governed
God's handmaids, nor his brother's daughters---not from suspicion, he
explained, but because the women who would necessarily accompany them
might place ``a stumbling-block or an occasion to fall\ldots\ in the way of
the weak'' (\work{Vita} 26). Preaching was the core of it, forty years of
it, with stenographers turning the spoken word into books as fast as he
made it; and behind the preaching the correspondence---to officials,
ladies, monks, pagans, Jerome in Bethlehem (an exchange of famous
prickliness on both sides), and eventually the whole Latin world.
